% from main.tex
%%%%%%%%%%  scan idx 264  |  folio 251  %%%%%%%%%%
immersion régulière et d'un morphisme lisse) est un morphisme parfait.

\textbf{Notation 4.2.} Soit \(f:X\to Y\) un morphisme localement de type fini de schémas. On notera \(D(f)_{\parf}\), ou \(D(X)_{f\text{-}\parf}\), ou \(D(X)_{Y\text{-}\parf}\) s'il n'y a pas de confusion à redouter (resp. \(D(f)_{\uparf}\), \(D(X)_{f\text{-}\uparf}\), \(D(X)_{Y\text{-}\uparf}\)) la sous-catégorie pleine de \(D(X)\) formée des complexes parfaits (resp. d'amplitude parfaite finie) rel. à \(f\). D'après (I 5.3) et (1.4), \(D(f)_{\parf}\) et \(D(f)_{\uparf}\) sont des sous-catégories triangulées de \(D(X)\). Si \(X\) est quasi-compact, on a \(D(f)_{\parf}=D(f)_{\uparf}\).

\textbf{Proposition 4.3} (cf. I 5.8). Soient \(f:X\to Y\) un morphisme localement de type fini de schémas, \(E\in\Ob D(X)\), \([a,b]\) un intervalle de \(\mathbb Z\), \(x\) un point de \(X\), \(y=f(x)\). Les conditions suivantes sont équivalentes :
\begin{enumerate}[label=(\roman*)]
\item \(\parfamp_f(E)\subset[a,b]\) au voisinage de \(x\) :
\item \(E\) est, au voisinage de \(x\), \((a-1)\)-pseudo-cohérent rel. à \(f\) et acyclique en dehors de \([a,b]\), et l'on a \(\tamp_{\mathcal O_{Y,y}}(E_x)\subset[a,b]\).
\end{enumerate}

\emph{Preuve.} L'implication (i) \(\Longrightarrow\) (ii) est triviale, prouvons (ii) \(\Longrightarrow\) (i). La question étant locale au voisinage de \(x\), on se ramène aussitôt, d'après les définitions, au cas où \(f\) est lisse. Alors \(E\) est \((a-1)\)-pseudo-cohérent au voisinage de \(x\), donc, quitte à se localiser, on peut supposer \(E\) strictement \((a-1)\)-pseudo-cohérent (I 2.1) et nul en degré \(>b\), et quitte à tronquer (i.e. à remplacer \(E\) par le complexe \((0\to E^a/B^a\to E^{a+1}\to\cdots E^b\to 0)\)), on peut supposer que \(E^i\) est libre de type fini pour \(i>a\), \(E^a\) de présentation finie, et \(E^i=0\) pour \(i\notin[a,b]\). D'après (I 5.1.1), \(E^a_x\) est plat rel. à \(\mathcal O_{Y,y}\), donc (EGA IV 11.3.1) \(E^a\) est \(f\)-plat au voisinage de \(x\). Posons pour abréger \(E^a=F\), il s'agit de montrer que \(F\) est pseudo-cohérent (au voisinage de \(x\)). Quitte à plonger \(X\) dans un espace affine type sur \(Y\), on peut
